\chapter{Manuscript Structure and Conventions}
\label{ch:meta_structure}

To aid the review of this manuscript, particularly concerning the internal consistency and the role of computer-assisted proofs, we provide here a dependency graph and a summary of conventions.

\section{Logical Spine and Dependency Graph}

The proof of the main theorem (Existence and Mass Gap) relies on decomposing the coupling space $\beta \in (0, \infty)$ into three regimes. The dependencies are structured as follows:

\begin{enumerate}
    \item \textbf{Core Analytic Framework (Chapter 3 \& 4):}
    \begin{itemize}
        \item \textbf{Theorem \ref{thm:weak-coupling-uniqueness} (UV Stability):} Purely Analytic. Relies on Balaban's Renormalization Group. Establishes the existence of the infinite volume limit and asymptotic freedom.
        \item \textbf{Gap Comparison (Thermodynamics \texorpdfstring{$\to$}{->} Hamiltonian):} Purely Analytic. Connects the LSI constant of the Glauber dynamics to the spectral gap of the physical Hamiltonian.
        \item \textbf{Theorem \ref{thm:stability_gap_fixed} (Continuum Limit):} Purely Analytic. Proves that Norm Resolvent Convergence preserves the spectral gap.
    \end{itemize}

    \item \textbf{Regime-Specific Inputs:}
    \begin{itemize}
    \item \textbf{Strong Coupling ($\beta \le \beta_{\mathrm{clust}}$):} Purely Analytic. Cluster Expansion (convergent polymer expansion), with $\beta_{\mathrm{clust}}\approx 0.016$.
        \item \textbf{Weak Coupling ($\beta > \beta_W$):} Purely Analytic. Bakry-\'{E}mery criterion + perturbative RG flow control.
    \item \textbf{Transition Regime ($\beta_{\mathrm{clust}} < \beta \le \beta_W$):} \textbf{Dobrushin/FSC bridge + Computer-Assisted Crossover Contraction.}
        \begin{itemize}
            \item \textbf{Interval Bridge (CAP):} Uses rigorous Interval Arithmetic to verify the contraction of the RG map via Ab Initio Coefficient Tracing.
            \item \textbf{Certification:} The CAP allows certification of the stability of the Wilson coefficients (Plaquette, Rectangle, etc.). This, in principle, implies the non-perturbative stability of the full Yang-Mills effective action.
        \end{itemize}
    \end{itemize}
\end{enumerate}

\subsection{Role of Computer-Assisted Verification}
We clarify the status of the intermediate regime. 
Specifically:
\begin{itemize}
    \item The CAP is designed to be **rigorous** for the full theory. It implements the "Ab Initio System," which traces the evolution of specific Wilson coefficients using intervals derived from the Haar measure integration.
    \item **Current Status:** The verification currently utilizes \textbf{Ab Initio} bounds for input coefficients, replacing the previous "Proxy Models". Consequently, the intermediate-regime conclusions are \textbf{CAP-certified} \emph{under the standard computational hypotheses} for rigorous interval verification (correct execution of the verifier, correctness of the interval arithmetic implementation, and the stated floating-point/arbitrary-precision model); see Appendix A and Chapter 8.
\end{itemize}

\section{Conventions and Normalizations}

To ensure consistency across the chapters, we fix the following conventions.

\subsection{Group Geometry and Metrics}
For the gauge group $G=SU(N)$:
\begin{itemize}
    \item We use the standard bi-invariant metric normalized such that the shortest geodesic length of a maximal torus is $2\pi$.
    \item \textbf{Log-Sobolev Constant:} We adopt the normalization where the Ricci curvature of $SU(N)$ is positive. The standard LSI constant for the heat kernel on $SU(N)$ is taken as:
    \begin{equation}
        \rho_{SU(N)} = \frac{2}{N-1}
    \end{equation}
    (Note: Any deviation from this in earlier drafts is superseded by this definition).
\end{itemize}

\subsection{Hamiltonian and Gaps}
\begin{itemize}
    \item \textbf{Lattice Hamiltonian:} $H = -\frac{1}{a} \ln T$, where $T$ is the transfer matrix per time slice.
    \item \textbf{Gap Definition:} The spectral gap $\Delta(H)$ is the infimum of the spectrum of $H$ restricted to the subspace orthogonal to the vacuum.
    \item \textbf{Comparison Normalization:} The inequality $\gap(H) \ge \frac{1}{2}\gamma$ refers to the spectral gap $\gamma$ of the define continuous-time Glauber dynamics generator $L$, normalized such that single-site updates happen at rate 1.
\end{itemize}

\section{Proof Status Table}

\begin{table}[h]
\centering
\begin{tabular}{|p{0.3\textwidth}|p{0.2\textwidth}|p{0.4\textwidth}|}
\hline
\textbf{Theorem/Claim} & \textbf{Status} & \textbf{Note} \\ \hline
Existence of Infinite Volume Limit & Proved & Balaban (1985) + adaptations in Ch 3. \\ \hline
UV Stability (Weak Coupling) & Proved & Theorem \ref{thm:weak-coupling-uniqueness}. Analytic. \\ \hline
Intermediate Crossover & \textbf{Conditional} & Computer-assisted interval-bridge verification (depends on rigorous interval execution). \\ \hline
Strong Coupling Gap & Proved & Rigorous cluster expansion for $\beta \le \beta_{\mathrm{clust}}\approx 0.016$, extended to $\beta \le \VerBetaStrongMax$ via Dobrushin/FSC bridge; see Ch.~1 and App.~\ref{app:hard_analysis_complete}. \\ \hline
Gap Comparison ($\Delta(H) \ge \gamma/2$) & Proved & Established in the LSI/transfer-matrix bridge (uniform in volume, conditioned on RP). \\ \hline
Continuum Limit (Norm Resolvent) & Proved & Section \ref{sec:resolvent_convergence_fixed}. Relies on uniform stability bounds. \\ \hline
Mass Gap for SU($N\ge 5$) & Conditional & Proof for Modified Action requires RP check \& Tuning. \\ \hline
Physical Scale Setting & Conditional & Defined via Gradient Flow; stability conditional on CAP. \\ \hline
\end{tabular}
\caption{Status of key components of the proof.}
\label{tab:proof_status}
\end{table}
